%% Heading level is chosen by the caller (\section in the short version,
%% \subsection in the full appendix).
\providecommand{\genexheading}{\section}
\genexheading{Selected Generations}
\label{app:generated-examples}

Each item gives the final compiled Lean statement followed by a mathematical
informalization. 

\begin{list}{}{%
  \setlength{\leftmargin}{0pt}%
  \setlength{\rightmargin}{0pt}%
  \setlength{\labelwidth}{0pt}%
  \setlength{\labelsep}{0pt}%
  \setlength{\itemindent}{0pt}%
  \setlength{\itemsep}{1.25em}%
  \setlength{\parsep}{0pt}%
  \setlength{\topsep}{0.75em}%
}

\item
\begin{samepage}
\begin{lstlisting}[basicstyle=\ttfamily\footnotesize,frame=single,breaklines=true]
∀ q : ℍ,
  q * q = -1 ↔ q.re = 0 ∧ ‖q‖ = 1
\end{lstlisting}
\noindent\textbf{Quaternionic square roots of minus one.}
A real quaternion squares to $-1$ if and only if its real component is zero
and its norm is one. Equivalently, the square roots of $-1$ are precisely the
unit purely imaginary quaternions.
\end{samepage}

\begin{comment}\item
\begin{samepage}
\begin{lstlisting}[basicstyle=\ttfamily\footnotesize,frame=single,breaklines=true]
∀ (α : Type _) [Group α], ∀ a b : α,
  a * b = 1 ↔
    b * a = 1 ∧
      ∀ c : α, c * a = 1 ↔ a * c = 1
\end{lstlisting}
\noindent\textbf{One-sided inverses in a group.}
For elements $a,b$ of a group, $ab=1$ exactly when $ba=1$ and, for every
$c$, the equation $ca=1$ is equivalent to $ac=1$. All of these equations say
that the other element is the unique inverse of $a$.
\end{samepage}
\end{comment}

\begin{comment}
\item
\begin{samepage}
\begin{lstlisting}[basicstyle=\ttfamily\footnotesize,frame=single,breaklines=true]
∀ x : ℝ,
  0 < x ↔
    Real.log (Real.exp x + 1) > Real.log 2
\end{lstlisting}
\noindent\textbf{Positivity through a logarithmic transform.}
A real number $x$ is positive if and only if
$\log(1+e^x)>\log 2$. Thus the softplus transform crosses its value at zero
exactly when $x$ crosses zero.
\end{samepage}
\end{comment}

\item
\begin{samepage}
\begin{lstlisting}[basicstyle=\ttfamily\footnotesize,frame=single,breaklines=true]
∀ z : UpperHalfPlane,
  ∃ n : ℕ, 0 < n ∧
    ‖Complex.exp
      (2 * ↑Real.pi * Complex.I * ↑z) ^ n‖ < 1 / 2
\end{lstlisting}
\noindent\textbf{Decay of the upper-half-plane exponential.}
For every $z$ in the complex upper half-plane, some positive power of
$e^{2\pi i z}$ has absolute value less than $1/2$. 
\end{samepage}

\begin{comment}
\item
\begin{samepage}
\begin{lstlisting}[basicstyle=\ttfamily\footnotesize,frame=single,breaklines=true]
∀ r : ℝ, r < 1 →
  (0 < r ∧ r < 1 ↔
    Real.log ((1 + r) / (1 - r)) > 0)
\end{lstlisting}
\noindent\textbf{Positivity of a logarithmic ratio.}
For a real $r<1$, the quantity
$\log\!\left((1+r)/(1-r)\right)$ is positive exactly when $0<r<1$.
\end{samepage}
\end{comment}

\item
\begin{samepage}
\begin{lstlisting}[basicstyle=\ttfamily\footnotesize,frame=single,breaklines=true]
∀ (n : Type _) [Fintype n] [DecidableEq n]
    (R : Type _) [CommRing R]
    (A : Matrix n n R),
  A * A = 1 ↔
    A.det * A.det = 1 ∧
      A.adjugate = A.det • A
\end{lstlisting}
\noindent\textbf{Involutory matrices and their adjugates.}
For a finite square matrix $A$ over a commutative ring, $A^2=I$ if and only if
$(\det A)^2=1$ and $\operatorname{adj}(A)=(\det A)A$.
\end{samepage}

\item
\begin{samepage}
\begin{lstlisting}[basicstyle=\ttfamily\footnotesize,frame=single,breaklines=true]
∀ κ : Cardinal,
  κ * κ = κ ↔
    κ = 0 ∨ κ = 1 ∨ Cardinal.aleph0 ≤ κ
\end{lstlisting}
\noindent\textbf{Multiplicatively idempotent cardinals.}
A cardinal is unchanged by squaring exactly when it is $0$, $1$, or infinite
(equivalently, at least $\aleph_0$).
\end{samepage}

\item
\begin{samepage}
\begin{lstlisting}[basicstyle=\ttfamily\footnotesize,frame=single,breaklines=true]
∀ z : UpperHalfPlane,
  ∃ w : UpperHalfPlane,
    w = -(1 : ℂ) / (z : ℂ) ∧
      ∀ n : ℕ, n > 0 ↔
        (↑n : ℂ) * (z : ℂ) ≠ 0
\end{lstlisting}
\noindent\textbf{Modular inversion in the upper half-plane.}
For every upper-half-plane point $z$, the modular inverse $w=-1/z$ is again in
the upper half-plane. Moreover, for every natural number $n$, positivity of
$n$ is equivalent to the nonvanishing of the complex product $nz$.
\end{samepage}

\item
\begin{samepage}
\begin{lstlisting}[basicstyle=\ttfamily\footnotesize,frame=single,breaklines=true]
∀ (R : Type _) [CommRing R],
  ∀ (W' : WeierstrassCurve.Jacobian R)
      (P : Fin 3 → R),
    W'.Equation (W'.neg P) ↔ W'.Equation P
\end{lstlisting}
\noindent\textbf{Negation preserves the Jacobian equation.}
For a Jacobian Weierstrass model over any commutative ring, negating a
projective point preserves the curve equation: $-P$ lies on the curve exactly
when $P$ does.
\end{samepage}

\item
\begin{samepage}
\begin{lstlisting}[basicstyle=\ttfamily\footnotesize,frame=single,breaklines=true]
∀ (K : Type _) [Field K] (f : Polynomial K),
  Function.Injective
    (fun g : Polynomial K => g * f) ↔
      f ≠ 0
\end{lstlisting}
\noindent\textbf{Injectivity of polynomial multiplication.}
Over a field, right multiplication by a polynomial $f$ is injective on the
polynomial ring if and only if $f$ is nonzero.
\end{samepage}

\item
\begin{samepage}
\begin{lstlisting}[basicstyle=\ttfamily\footnotesize,frame=single,breaklines=true]
∀ (K : Type _) [Field K] [NumberField K]
    (x : K),
  x ≠ 0 ↔
    ∀ w : NumberField.InfinitePlace K,
      (NumberField.InfinitePlace.mk
        (NumberField.InfinitePlace.embedding w)) x > 0
\end{lstlisting}
\noindent\textbf{Positivity at every infinite place.}
An element of a number field is nonzero exactly when its value at every
infinite place is strictly positive, where an infinite place evaluates through
the absolute value associated with its embedding.
\end{samepage}

\item
\begin{samepage}
\begin{lstlisting}[basicstyle=\ttfamily\footnotesize,frame=single,breaklines=true]
∀ (α : Type _) [GeneralizedBooleanAlgebra α]
    (a b : α),
  a = b ↔
    ∀ c : α, a \ c = b \ c ∧ c \ a = c \ b
\end{lstlisting}
\noindent\textbf{Equality via relative complements.}
Two elements $a,b$ of a generalized Boolean algebra are equal exactly when,
for every $c$, removing $c$ from them gives the same result and removing them
from $c$ also gives the same result.
\end{samepage}

\item
\begin{samepage}
\begin{lstlisting}[basicstyle=\ttfamily\footnotesize,frame=single,breaklines=true]
∀ (V : Type _) [Finite V] (G : SimpleGraph V),
  G.edgeSet = ∅ ↔
    ∀ s : Set V, G.IsClique s ↔ s.ncard < 2
\end{lstlisting}
\noindent\textbf{Empty graphs characterized by their cliques.}
A finite simple graph has no edges if and only if its cliques are precisely the
vertex sets containing fewer than two vertices.
\end{samepage}

\item
\begin{samepage}
\begin{lstlisting}[basicstyle=\ttfamily\footnotesize,frame=single,breaklines=true]
∀ (V : Type _) (G : SimpleGraph V) [Finite V],
  G.IsClique (Set.univ : Set V) ↔
    ∀ (M : Gᶜ.Subgraph),
      M.IsMatching ↔ M.verts = ∅
\end{lstlisting}
\noindent\textbf{Completeness through matchings in the complement.}
A finite simple graph is complete if and only if a subgraph of its complement
is a matching exactly when that subgraph has no vertices.
\end{samepage}

\item
\begin{samepage}
\begin{lstlisting}[basicstyle=\ttfamily\footnotesize,frame=single,breaklines=true]
∀ n : ℕ, n > 0 ↔
  ∀ f : Polynomial ℝ, f.natDegree ≤ n - 1 →
    ((∀ k < n,
      Polynomial.eval
        (Polynomial.Chebyshev.node n k) f = 0) ↔
      f = 0)
\end{lstlisting}
\noindent\textbf{Vanishing at Chebyshev nodes.}
A natural number $n$ is positive exactly when every real polynomial of degree
at most $n-1$ that vanishes at all $n$ Chebyshev nodes is the zero polynomial
(and the zero polynomial vanishes at those nodes).
\end{samepage}

\item
\begin{samepage}
\begin{lstlisting}[basicstyle=\ttfamily\footnotesize,frame=single,breaklines=true]
∀ (C : Type _) [CategoryTheory.Category C]
    [CategoryTheory.Abelian C]
    (X Y : C) (f : X ⟶ Y),
  CategoryTheory.Mono f →
    (CategoryTheory.IsIso f ↔
      ∀ (Z : C) (g : Y ⟶ Z),
        (g = 0 ↔ f ≫ g = 0))
\end{lstlisting}
\noindent\textbf{Detecting isomorphisms in an abelian category.}
Let $f:X\to Y$ be a monomorphism in an abelian category. Then $f$ is an
isomorphism exactly when precomposition with $f$ detects zero morphisms out of
$Y$: for every $g:Y\to Z$, one has $g=0$ if and only if $g\circ f=0$.
\end{samepage}

\end{list}
